\chapter{Introducción}
\label{sec:introduccion}

El presente informe detalla el desarrollo, la arquitectura de firmware y la aplicación de control para la central
electrónica de semáforos diseñada en el marco del Hackathon 2026. El propósito del proyecto consistió en concebir e
implementar un sistema capaz de regular el tránsito vehicular y peatonal en una intersección de dos fases
ortogonales (Norte--Sur y Este--Oeste).

El núcleo de control fue desarrollado sobre la plataforma embebida STM32F103 (Blue Pill), cuyo microcontrolador
ARM Cortex-M3 opera junto con el sistema operativo de tiempo real FreeRTOS. A partir de una arquitectura multitarea
desacoplada, el firmware gestiona la máquina de estados de las señales luminosas, ejecuta la adquisición periódica de
sensores y variables de diagnóstico, y atiende peticiones de comunicación serie industrial bajo el protocolo Modbus RTU
en modalidad esclavo.

De manera complementaria, se diseñó (con el uso de la IA) una aplicación de supervisión y comando que asume el rol de
maestro en la red.  Esta herramienta permite auditar en tiempo real la telemetría del controlador, conmutar modos de
operación en caliente y reprogramar los intervalos temporales de cada fase sin interrumpir la ejecución del sistema. En
los capítulos siguientes se expone detalladamente el funcionamiento interno del firmware, el esquema de estados y la
integración con la interfaz de usuario.
